\section{Feature 005: Add a Choice}
\label{sec:feature005}

\subsection{Mô tả và hướng dẫn sử dụng}
Quản trị viên tạo hoạt động Choice từ trang khóa học F005 (URL: \texttt{\textless{}BASE\_URL\textgreater{}/course/view.php?id=2}) qua \textbf{Edit mode $>$ General $>$ Add an activity or resource $>$ Choice $>$ Add}. Form gồm các trường \textit{Choice name}, \textit{Description}, \textit{Display mode for the options} và các nhóm \textit{Options}, \textit{Availability}, \textit{Results}. \textit{Choice name} và \textit{Option 1} là hai trường bắt buộc. Các nhóm thiết lập chung của hoạt động như \textit{Common module settings}, \textit{Restrict access} và \textit{Activity completion} được giữ ở giá trị mặc định trong bộ kiểm thử.

\textbf{Các bước sử dụng:}
\begin{enumerate}[noitemsep]
  \item Đăng nhập bằng tài khoản quản trị.
  \item Truy cập khóa học và bật \textit{Edit mode}. Trong mục \textit{General}, bấm \textit{Add an activity or resource}, chọn \textit{Choice} rồi bấm \textit{Add}.
  \item Nhập hai trường bắt buộc là \textit{Choice name} và \textit{Option 1}. Điền các phương án bổ sung, mô tả, giới hạn phản hồi, thời gian và chế độ công bố kết quả theo yêu cầu. Giữ mặc định các thiết lập chung của hoạt động.
  \item Bấm \textbf{Save and display} để lưu và mở trang Choice, hoặc \textbf{Save and return to course} để lưu và quay về trang khóa học. Bấm \textbf{Cancel} để hủy tạo.
  \item Khi lưu, Moodle kiểm tra dữ liệu. Nếu dữ liệu hợp lệ, hoạt động được tạo và Moodle mở trang tương ứng với nút đã chọn. Nếu thiếu dữ liệu bắt buộc hoặc giờ đóng sớm hơn giờ mở, form vẫn mở và hiển thị lỗi tại trường tương ứng để quản trị viên sửa rồi lưu lại.
\end{enumerate}
Trong bộ test F005, người kiểm thử tắt \textit{Edit mode} sau khi kiểm tra kết quả để chuẩn bị cho ca tiếp theo.

\subsection{Use case và activity diagram}
\begin{table}[H]
\centering\small
\begin{tabular}{|L{3.6cm}|L{11.4cm}|}
\hline
Tên & UC-005: Tạo hoạt động Choice \\ \hline
Tác nhân & Quản trị viên (Admin) \\ \hline
Tiền điều kiện & Quản trị viên đã đăng nhập. Khóa thử nghiệm có mục General và cho phép thêm hoạt động Choice. Edit mode đang tắt. \\ \hline
Hậu điều kiện & Choice được lưu trong khóa học. Cuối ca, Edit mode được tắt. \\ \hline
Luồng chính & Mở form, nhập tên và hai phương án, chọn Save and display rồi kiểm tiêu đề Choice. \\ \hline
Luồng thay thế & A1: lưu và về khóa học. A2: thêm ô phương án. A3: hủy tạo. \\ \hline
Luồng ngoại lệ & E1: thiếu tên. E2: thiếu Option 1. E3: giờ đóng trước giờ mở. Moodle báo lỗi tại trường để admin sửa và lưu lại. \\ \hline
\end{tabular}
\caption{Đặc tả use case UC-005}
\end{table}

\begin{figure}[H]
\centering
\begin{tikzpicture}[>=Stealth,
act/.style={draw,rounded corners,align=center,inner sep=5pt,fill=blue!5},
dec/.style={diamond,draw,aspect=2,align=center,inner sep=1pt,fill=yellow!15}]
\node[act] (open) {Mở form Choice};
\node[act,below=of open] (fill) {Nhập dữ liệu};
\node[dec,below=of fill] (cancel) {Hủy?};
\node[act,right=2cm of cancel] (back) {Về khóa học};
\node[dec,below=of cancel] (valid) {Lưu hợp lệ?};
\node[act,left=1.7cm of valid] (error) {Hiện lỗi tại trường};
\node[act,below=of valid] (saved) {Hiện Choice hoặc trang khóa học};
\node[act,below=of saved] (end) {Kiểm kết quả và tắt Edit mode};
\draw[->] (open)--(fill);
\draw[->] (fill)--(cancel);
\draw[->] (cancel)--node[above]{Có}(back);
\draw[->] (cancel)--node[right]{Không}(valid);
\draw[->] (valid)--node[above]{Không}(error);
\draw[->] (error)|-(fill);
\draw[->] (valid)--node[right]{Có}(saved);
\draw[->] (saved)--(end);
\draw[->] (back)|-(end);
\end{tikzpicture}
\caption{Luồng tạo, sửa lỗi và hủy hoạt động Choice}
\end{figure}

\subsection{Các biến đầu vào}
\begin{table}[H]
\centering\small
\begin{tabular}{|L{3.2cm}|L{2.1cm}|L{4.4cm}|L{4.3cm}|}
\hline
\rowcolor{gray!15}\textbf{Biến} & \textbf{Kiểu} & \textbf{Dữ liệu} & \textbf{Ràng buộc} \\ \hline
Choice name & Chuỗi & Tên hoạt động & Bắt buộc, tối đa 1.333 ký tự \\ \hline
Description và hiển thị & Văn bản, cờ & Mô tả, hiển thị trên khóa, bố cục & Phụ thuộc cấu hình Moodle \\ \hline
Options & Chuỗi, cờ & Phương án, đổi lựa chọn, chọn nhiều & Option 1 bắt buộc \\ \hline
Limit & Số, cờ & Bật giới hạn, số lượt, hiện chỗ còn & Limit chỉ hiện khi bật giới hạn \\ \hline
Availability & Ngày giờ, cờ & Giờ mở, giờ đóng, xem trước & Khi bật cả hai, giờ đóng không trước giờ mở \\ \hline
Results & Danh sách, cờ & Công bố, quyền riêng tư, chưa trả lời, tài khoản không hoạt động & Quyền riêng tư phụ thuộc chế độ công bố \\ \hline
\end{tabular}
\caption{Các biến đầu vào của Feature 005}
\end{table}

\subsection{Kỹ thuật 1: Phân lớp tương đương (ECP)}
Nhóm phân lớp dữ liệu theo yêu cầu nhập liệu và sự phụ thuộc giữa các tùy chọn trên form. Mỗi ca đại diện kiểm tra thông báo tại trường hoặc giá trị còn hiển thị sau khi lưu. Các mã rút gọn trong bảng là ba chữ số cuối của mã ca kiểm thử.

\begin{longtable}{|L{4.8cm}|L{9.5cm}|}
\hline
\textbf{Đầu vào} & \textbf{Phân lớp và ca đại diện}\\ \hline
Choice name & Tên thường hợp lệ (001), tên rỗng (002), độ dài ở biên (003--006).\\ \hline
Nội dung và số ô phương án & Option 1 trống (007), chỉ có một phương án (008), hai phương án (001), thêm ô để nhập phương án thứ sáu (009).\\ \hline
Description và công tắc hiển thị & Để mặc định (001), nhập mô tả và bật hiển thị (010).\\ \hline
Display mode, Allow update, Allow multiple & So sánh thiết lập mặc định ở 001 với hiển thị dọc ở 011, cho đổi câu trả lời ở 012 và chọn nhiều ở 013.\\ \hline
Limit responses, Limit và Show available spaces & Tắt giới hạn (014), bật giới hạn và nhập 0, 1, 2 (015--017). Ca 017 bật hiển thị số chỗ còn.\\ \hline
Giờ mở, giờ đóng và Show preview & Tắt cả hai (001), chỉ bật giờ mở và xem trước (018), chỉ bật giờ đóng (019), bật cả hai với quan hệ thời gian khác nhau (020--021).\\ \hline
Publish results, Privacy of results, Show unanswered & Công bố theo mặc định (001), bật cột chưa trả lời (022), luôn công bố ẩn danh (023). Các biến thể công bố, quyền riêng tư và Include inactive ngoài các ca này chưa có ca riêng.\\ \hline
Thao tác cuối & Lưu và mở Choice (001), lưu và trở về khóa (010), hủy (024).\\ \hline
\end{longtable}

\subsection{Kỹ thuật 2: Phân tích giá trị biên (BVA)}
\begin{itemize}[noitemsep]
\item Tên Choice: 0, 1, 2 ký tự (002--004), 1.332 và 1.333 ký tự (005--006). Giá trị 1.334 ký tự nằm ngoài bộ 24 ca hiện tại.
\item Số phương án hợp lệ: ca 007 kiểm Option 1 bắt buộc dù đã nhập Option 2, ca 008 kiểm một phương án và ca 001 kiểm hai phương án. Ca 009 kiểm việc bổ sung ô sau năm ô ban đầu.
\item Giới hạn phản hồi: kiểm 0, 1 và 2 ở ca 015--017. Kết quả của các ca này là giá trị đã lưu trên form Settings.
\item Khoảng cách giờ đóng so với giờ mở: ca 020 thử $-1$ phút, kiểm lỗi rồi sửa thành $+1$ phút và lưu lại. Ca 021 thử đúng 0 phút và kiểm cấu hình được lưu.
\end{itemize}

\subsection{Kỹ thuật 3: Bảng quyết định (Decision table)}
\begin{longtable}{|L{5.7cm}|L{6.9cm}|L{1.7cm}|}
\hline
\textbf{Điều kiện} & \textbf{Kết quả mong đợi} & \textbf{Ca}\\ \hline
Tên trống hoặc Option 1 trống & Báo lỗi ở trường tương ứng. Sửa dữ liệu rồi lưu được. & 002, 007\\ \hline
Giới hạn tắt & Các ô Limit ẩn. & 014\\ \hline
Giới hạn bật, nhập 0/1/2 & Ô Limit hiển thị và lưu giá trị. Khi bật Show available spaces, lưu thiết lập Yes. & 015--017\\ \hline
Chỉ bật giờ mở hoặc chỉ bật giờ đóng & Lưu mốc đã bật. Ca 018 lưu thêm Show preview. & 018--019\\ \hline
Bật cả hai, đóng trước giờ mở & Báo lỗi ở giờ đóng. Sửa thành giờ sau lúc mở rồi lưu được. & 020\\ \hline
Bật cả hai, đóng bằng giờ mở & Lưu hai mốc bằng nhau. & 021\\ \hline
Luôn công bố và chọn ẩn danh & Lưu hai thiết lập công bố và quyền riêng tư. & 023\\ \hline
\end{longtable}

\subsection{Kỹ thuật 4: Use-case testing}
\begin{table}[H]
\centering\small
\begin{tabular}{|c|L{7.0cm}|L{3.8cm}|c|}
\hline
\rowcolor{gray!15}\textbf{Luồng} & \textbf{Kịch bản} & \textbf{Kết quả mong đợi} & \textbf{Ca} \\ \hline
P1 & Lưu và hiển thị Choice & Mở trang Choice & 001 \\ \hline
A1 & Lưu và quay về khóa học & Hiện Choice và mô tả & 010 \\ \hline
A2 & Thêm ô phương án rồi lưu & Lưu phương án bổ sung & 009 \\ \hline
A3 & Hủy tạo Choice & Không tạo hoạt động & 024 \\ \hline
E1 & Thiếu tên, nhập lại rồi lưu & Báo lỗi rồi lưu được & 002 \\ \hline
E2 & Thiếu Option 1, nhập lại rồi lưu & Báo lỗi rồi lưu được & 007 \\ \hline
E3 & Giờ đóng trước giờ mở, sửa rồi lưu & Báo lỗi rồi lưu được & 020 \\ \hline
\end{tabular}
\caption{Scenario và test case use-case testing của Feature 005}
\end{table}

\subsection{Kết quả thực thi và phát hiện}\label{sec:f005-result}
Ở ca TC-005-001, Moodle hiển thị hoạt động \texttt{F005-CHOICE-001} cùng hai phương án đã nhập. Katalon kiểm tra được tiêu đề, nhưng kịch bản ghi ban đầu có một bước mở trực tiếp trang Choice. Vì vậy, lượt chạy này chưa xác nhận được thao tác \textit{Save and display} tự chuyển đến hoạt động vừa tạo. Kết quả thực thi của cả bộ 24 ca chưa được tổng hợp.

\subsection{Danh sách test case của Feature 005}
Bảng dưới đây tổng hợp 24 ca kiểm thử của Feature 005. Các bước thao tác và kết quả mong đợi được ghi trong bảng test case của nhóm.
{\footnotesize
\begin{longtable}{|l|L{4.4cm}|c|L{6.4cm}|}
\hline
\rowcolor{gray!15}\textbf{ID} & \textbf{Tên test case} & \textbf{Kỹ thuật} & \textbf{Mô tả} \\ \hline
\endhead
TC-005-001 & Tạo Choice cơ bản & ECP/UC & Tạo Choice cơ bản. \\ \hline
TC-005-002 & Thiếu tên Choice & ECP/BVA & Thiếu tên Choice. Sửa trường bị lỗi và lưu lại thành công. \\ \hline
TC-005-003 & Tên Choice 1 ký tự & BVA & Tên Choice 1 ký tự. \\ \hline
TC-005-004 & Tên Choice 2 ký tự & BVA & Tên Choice 2 ký tự. \\ \hline
TC-005-005 & Tên Choice 1.332 ký tự & BVA & Tên Choice 1.332 ký tự. \\ \hline
TC-005-006 & Tên Choice 1.333 ký tự & BVA & Tên Choice 1.333 ký tự. \\ \hline
TC-005-007 & Thiếu Option 1 & ECP/BVA & Thiếu Option 1. Sửa trường bị lỗi và lưu lại thành công. \\ \hline
TC-005-008 & Chỉ nhập Option 1 & ECP/BVA & Chỉ nhập Option 1. \\ \hline
TC-005-009 & Thêm Option 6 & ECP/UC & Thêm Option 6. \\ \hline
TC-005-010 & Hiện mô tả trên trang khóa học & ECP/DT & Hiện mô tả trên trang khóa học. \\ \hline
TC-005-011 & Hiển thị theo chiều dọc & ECP & Hiển thị theo chiều dọc. \\ \hline
TC-005-012 & Cho phép đổi câu trả lời & ECP & Cho phép đổi câu trả lời. \\ \hline
TC-005-013 & Cho phép chọn nhiều phương án & ECP & Cho phép chọn nhiều phương án. \\ \hline
TC-005-014 & Tắt giới hạn lượt chọn & ECP/DT & Tắt giới hạn lượt chọn. \\ \hline
TC-005-015 & Đặt Limit 1 bằng 0 & BVA/DT & Đặt Limit 1 bằng 0. \\ \hline
TC-005-016 & Đặt Limit 1 bằng 1 & BVA/DT & Đặt Limit 1 bằng 1. \\ \hline
TC-005-017 & Đặt Limit 1 bằng 2 và hiện chỗ còn & BVA/DT & Đặt Limit 1 bằng 2 và hiện chỗ còn. \\ \hline
TC-005-018 & Chỉ bật giờ mở và bật xem trước & ECP/DT & Chỉ bật giờ mở và bật xem trước. \\ \hline
TC-005-019 & Chỉ bật giờ đóng & ECP/DT & Chỉ bật giờ đóng. \\ \hline
TC-005-020 & Giờ đóng trước giờ mở 1 phút & BVA/DT & Giờ đóng trước giờ mở 1 phút. Sửa trường bị lỗi và lưu lại thành công. \\ \hline
TC-005-021 & Giờ đóng bằng giờ mở & BVA/DT & Giờ đóng bằng giờ mở. \\ \hline
TC-005-022 & Hiện cột người chưa trả lời & ECP & Bật Show column for unanswered và kiểm cài đặt được lưu. \\ \hline
TC-005-023 & Công bố kết quả ẩn danh & ECP/DT & Công bố kết quả ẩn danh. \\ \hline
TC-005-024 & Hủy tạo Choice & UC & Hủy tạo Choice. \\ \hline
\end{longtable}
}


\subsection{Test case mẫu của Feature 005}
Các ca mẫu trình bày luồng tạo thành công, cách xử lý khi thiếu dữ liệu bắt buộc, thiết lập giới hạn phản hồi, sửa giờ đóng và hủy tạo hoạt động.
\bigskip\noindent\textbf{TC-005-001 -- Tạo Choice cơ bản}\\
\noindent\textit{Precondition:} Moodle demo \url{http://127.0.0.1:9501/}, giao diện English. Đã đăng nhập admin. Khóa F005-ENROL-LAB có mục General và bật hoạt động Choice. Chưa có Choice dùng dữ liệu tên của ca này. Edit mode ở trạng thái Off trước khi chạy.
\smallskip

{\footnotesize
\begin{longtable}{|c|L{3.0cm}|L{5.1cm}|L{5.6cm}|}
\hline
\rowcolor{gray!15}\textbf{\#} & \textbf{Bước} & \textbf{Mô tả} & \textbf{Kết quả mong đợi} \\ \hline
\endhead
1 & Mở khóa học & Vào Home  $\rightarrow$  F005-ENROL-LAB. & Trang khóa học hiển thị. \\ \hline
2 & Bật Edit mode & Bật Edit mode. & Hiện nút thêm hoạt động trong khóa. \\ \hline
3 & Mở form Choice & Trong General, bấm Add an activity or resource, chọn Choice rồi bấm Add. & Form tạo Choice hiển thị. \\ \hline
4 & Nhập Choice name & Nhập F005-CHOICE-001 vào Choice name. & Choice name có 15 ký tự. \\ \hline
5 & Nhập Option 1 & Nhập Sáng vào Option 1. & Option 1 hiện Sáng. \\ \hline
6 & Nhập Option 2 & Nhập Chiều vào Option 2. & Option 2 hiện Chiều. \\ \hline
7 & Bấm lưu & Bấm Save and display. & Trang Choice mở với tên F005-CHOICE-001. \\ \hline
8 & Tắt Edit mode & Bấm Home, mở lại khóa F005 rồi tắt Edit mode. & Edit mode hiển thị Off. \\ \hline
\end{longtable}
}

\bigskip\noindent\textbf{TC-005-002 -- Thiếu tên Choice}\\
\noindent\textit{Precondition:} Moodle demo \url{http://127.0.0.1:9501/}, giao diện English. Đã đăng nhập admin. Khóa F005-ENROL-LAB có mục General và bật hoạt động Choice. Chưa có Choice dùng dữ liệu tên của ca này. Edit mode ở trạng thái Off trước khi chạy.
\smallskip

{\footnotesize
\begin{longtable}{|c|L{3.0cm}|L{5.1cm}|L{5.6cm}|}
\hline
\rowcolor{gray!15}\textbf{\#} & \textbf{Bước} & \textbf{Mô tả} & \textbf{Kết quả mong đợi} \\ \hline
\endhead
1 & Mở khóa học & Vào Home  $\rightarrow$  F005-ENROL-LAB. & Trang khóa học hiển thị. \\ \hline
2 & Bật Edit mode & Bật Edit mode. & Hiện nút thêm hoạt động trong khóa. \\ \hline
3 & Mở form Choice & Trong General, bấm Add an activity or resource, chọn Choice rồi bấm Add. & Form tạo Choice hiển thị. \\ \hline
4 & Để trống Choice name & Để trống Choice name. & Choice name trống. \\ \hline
5 & Nhập Option 1 & Nhập Sáng vào Option 1. & Option 1 hiện Sáng. \\ \hline
6 & Nhập Option 2 & Nhập Chiều vào Option 2. & Option 2 hiện Chiều. \\ \hline
7 & Bấm lưu & Bấm Save and display. & Form vẫn mở. Choice name báo lỗi bắt buộc. \\ \hline
8 & Nhập lại Choice name & Nhập F005-CHOICE-002 vào Choice name. & Choice name hiện F005-CHOICE-002. \\ \hline
9 & Bấm lưu lại & Bấm Save and display. & Trang Choice hiển thị F005-CHOICE-002. \\ \hline
10 & Tắt Edit mode & Bấm Home, mở lại khóa F005 rồi tắt Edit mode. & Edit mode hiển thị Off. \\ \hline
\end{longtable}
}

\bigskip\noindent\textbf{TC-005-007 -- Thiếu Option 1}\\
\noindent\textit{Precondition:} Moodle demo \url{http://127.0.0.1:9501/}, giao diện English. Đã đăng nhập admin. Khóa F005-ENROL-LAB có mục General và bật hoạt động Choice. Chưa có Choice dùng dữ liệu tên của ca này. Edit mode ở trạng thái Off trước khi chạy.
\smallskip

{\footnotesize
\begin{longtable}{|c|L{3.0cm}|L{5.1cm}|L{5.6cm}|}
\hline
\rowcolor{gray!15}\textbf{\#} & \textbf{Bước} & \textbf{Mô tả} & \textbf{Kết quả mong đợi} \\ \hline
\endhead
1 & Mở khóa học & Vào Home  $\rightarrow$  F005-ENROL-LAB. & Trang khóa học hiển thị. \\ \hline
2 & Bật Edit mode & Bật Edit mode. & Hiện nút thêm hoạt động trong khóa. \\ \hline
3 & Mở form Choice & Trong General, bấm Add an activity or resource, chọn Choice rồi bấm Add. & Form tạo Choice hiển thị. \\ \hline
4 & Nhập Choice name & Nhập F005-CHOICE-008 vào Choice name. & Choice name có 15 ký tự. \\ \hline
5 & Để trống Option 1 & Để trống Option 1. & Option 1 trống. \\ \hline
6 & Nhập Option 2 & Nhập Chiều vào Option 2. & Option 2 hiện Chiều. \\ \hline
7 & Bấm lưu & Bấm Save and display. & Form vẫn mở. Option 1 báo phải nhập ít nhất một phương án. \\ \hline
8 & Nhập lại Option 1 & Nhập Sáng vào Option 1. & Option 1 hiện Sáng. \\ \hline
9 & Bấm lưu lại & Bấm Save and display. & Trang Choice hiển thị F005-CHOICE-008. \\ \hline
10 & Tắt Edit mode & Bấm Home, mở lại khóa F005 rồi tắt Edit mode. & Edit mode hiển thị Off. \\ \hline
\end{longtable}
}

\bigskip\noindent\textbf{TC-005-017 -- Đặt Limit 1 bằng 2 và hiện chỗ còn}\\
\noindent\textit{Precondition:} Moodle demo \url{http://127.0.0.1:9501/}, giao diện English. Đã đăng nhập admin. Khóa F005-ENROL-LAB có mục General và bật hoạt động Choice. Chưa có Choice dùng dữ liệu tên của ca này. Edit mode ở trạng thái Off trước khi chạy.
\smallskip

{\footnotesize
\begin{longtable}{|c|L{3.0cm}|L{5.1cm}|L{5.6cm}|}
\hline
\rowcolor{gray!15}\textbf{\#} & \textbf{Bước} & \textbf{Mô tả} & \textbf{Kết quả mong đợi} \\ \hline
\endhead
1 & Mở khóa học & Vào Home  $\rightarrow$  F005 & Trang khóa học hiển thị. \\ \hline
2 & Bật Edit mode & Bật Edit mode. & Hiện nút thêm hoạt động trong khóa. \\ \hline
3 & Mở form Choice & Trong General, bấm Add an activity or resource, chọn Choice rồi bấm Add. & Form tạo Choice hiển thị. \\ \hline
4 & Nhập Choice name & Nhập F005-CHOICE-017 vào Choice name. & Choice name có 15 ký tự. \\ \hline
5 & Nhập Option 1 & Nhập Sáng vào Option 1. & Option 1 hiện Sáng. \\ \hline
6 & Nhập Option 2 & Nhập Chiều vào Option 2. & Option 2 hiện Chiều. \\ \hline
7 & Bật giới hạn & Chọn Yes ở Limit the number of responses allowed. & Hiện các ô Limit. \\ \hline
8 & Nhập Limit 1 & Nhập 2 vào Limit 1 và 2 vào Limit 2. & Limit 1 hiện 2. \\ \hline
9 & Hiện chỗ còn & Chọn Yes ở Show available spaces. & Hiện Yes. \\ \hline
10 & Bấm lưu & Bấm Save and display. & Trang Choice mở với tên F005-CHOICE-017. \\ \hline
11 & Kiểm cấu hình đã lưu & Bấm Settings của Choice. Đối chiếu trường đã chỉnh. & Limit 1 lưu 2, Limit 2 lưu 2. Show available spaces lưu Yes. \\ \hline
12 & Tắt Edit mode & Bấm Home, mở lại khóa F005 rồi tắt Edit mode. & Edit mode hiển thị Off. \\ \hline
\end{longtable}
}

\bigskip\noindent\textbf{TC-005-020 -- Giờ đóng trước giờ mở 1 phút}\\
\noindent\textit{Precondition:} Moodle demo \url{http://127.0.0.1:9501/}, giao diện English. Đã đăng nhập admin. Khóa F005-ENROL-LAB có mục General và bật hoạt động Choice. Chưa có Choice dùng dữ liệu tên của ca này. Edit mode ở trạng thái Off trước khi chạy.
\smallskip

{\footnotesize
\begin{longtable}{|c|L{3.0cm}|L{5.1cm}|L{5.6cm}|}
\hline
\rowcolor{gray!15}\textbf{\#} & \textbf{Bước} & \textbf{Mô tả} & \textbf{Kết quả mong đợi} \\ \hline
\endhead
1 & Mở khóa học & Vào Home  $\rightarrow$  F005 & Trang khóa học hiển thị. \\ \hline
2 & Bật Edit mode & Bật Edit mode. & Hiện nút thêm hoạt động trong khóa. \\ \hline
3 & Mở form Choice & Trong General, bấm Add an activity or resource, chọn Choice rồi bấm Add. & Form tạo Choice hiển thị. \\ \hline
4 & Nhập Choice name & Nhập F005-CHOICE-020 vào Choice name. & Choice name có 15 ký tự. \\ \hline
5 & Nhập Option 1 & Nhập Sáng vào Option 1. & Option 1 hiện Sáng. \\ \hline
6 & Nhập Option 2 & Nhập Chiều vào Option 2. & Option 2 hiện Chiều. \\ \hline
7 & Đặt giờ mở & Tích Enable tại Allow responses from. Chọn 01/12/2027 09:30. & Giờ mở hiện 01/12/2027 09:30. \\ \hline
8 & Đặt giờ đóng & Tích Enable tại Allow responses until. Chọn 01/12/2027 09:29. & Giờ đóng hiện 01/12/2027 09:29. \\ \hline
9 & Bấm lưu & Bấm Save and display. & Form vẫn mở. Giờ đóng báo lỗi vì trước giờ mở. \\ \hline
10 & Sửa giờ đóng & Trong Availability, đổi Allow responses until thành 01/12/2027 09:31. & Giờ đóng hiển thị 01/12/2027 09:31. \\ \hline
11 & Bấm lưu lại & Bấm Save and display. & Trang Choice hiển thị F005-CHOICE-020. \\ \hline
12 & Tắt Edit mode & Bấm Home, mở lại khóa F005 rồi tắt Edit mode. & Edit mode hiển thị Off. \\ \hline
\end{longtable}
}

\bigskip\noindent\textbf{TC-005-024 -- Hủy tạo Choice}\\
\noindent\textit{Precondition:} Moodle demo \url{http://127.0.0.1:9501/}, giao diện English. Đã đăng nhập admin. Khóa F005-ENROL-LAB có mục General và bật hoạt động Choice. Chưa có Choice dùng dữ liệu tên của ca này. Edit mode ở trạng thái Off trước khi chạy.
\smallskip

{\footnotesize
\begin{longtable}{|c|L{3.0cm}|L{5.1cm}|L{5.6cm}|}
\hline
\rowcolor{gray!15}\textbf{\#} & \textbf{Bước} & \textbf{Mô tả} & \textbf{Kết quả mong đợi} \\ \hline
\endhead
1 & Mở khóa học & Vào Home  $\rightarrow$  F005 & Trang khóa học hiển thị. \\ \hline
2 & Bật Edit mode & Bật Edit mode. & Hiện nút thêm hoạt động trong khóa. \\ \hline
3 & Mở form Choice & Trong General, bấm Add an activity or resource, chọn Choice rồi bấm Add. & Form tạo Choice hiển thị. \\ \hline
4 & Nhập Choice name & Nhập F005-CHOICE-024 vào Choice name. & Choice name có 15 ký tự. \\ \hline
5 & Nhập Option 1 & Nhập Sáng vào Option 1. & Option 1 hiện Sáng. \\ \hline
6 & Nhập Option 2 & Nhập Chiều vào Option 2. & Option 2 hiện Chiều. \\ \hline
7 & Bấm Cancel & Bấm Cancel ở cuối form. & Quay về trang khóa học. \\ \hline
8 & Kiểm danh sách & Vào Home  $\rightarrow$  F005 & Không có hoạt động F005-CHOICE-024. \\ \hline
9 & Tắt Edit mode & Bấm Home, mở lại khóa F005 rồi tắt Edit mode. & Edit mode hiển thị Off. \\ \hline
\end{longtable}
}

